\section{Результаты моделирования}

\subsection{Параметры модели}

Расчёт выполняется для радиального подшипника скольжения с геометрией~\varGeomLetter{} ($R = \varRmm$~мм, $c = \varCmkm$~мкм, $L = \varLmm$~мм) и \varDimpleTypeGenitive{} микрорельефом (тип~6 -- конический). Сводные исходные данные представлены в таблице~\ref{tab:input_data}.

\begin{table}[H]
\small
\caption{Исходные данные для расчёта}
\label{tab:input_data}
\begin{tabularx}{\textwidth}{|X|c|c|}
\hline
\centering\textbf{Параметр} & \textbf{Значение} & \textbf{Размерность} \\
\hline
Радиус вала $R$ & \varRmm & мм \\
\hline
Радиальный зазор $c$ & \varCmkm & мкм \\
\hline
Длина подшипника $L$ & \varLmm & мм \\
\hline
Относительная длина $L/D$ & \varLD & -- \\
\hline
Полуось по оси $Z$, $a$ & \varDimpleAmm & мм \\
\hline
Полуось по оси $\varphi$, $b$ & \varDimpleBmm & мм \\
\hline
Глубина углубления $h_p$ & \varDimpleHpmkm & мкм \\
\hline
Относительная глубина $h_p/c$ & \varDimpleHpRatio & -- \\
\hline
Число углублений по окружности $N_\varphi$ & \varNphi & шт. \\
\hline
Число рядов по длине $N_Z$ & \varNZ & шт. \\
\hline
Область текстурирования & \varPhiRange & -- \\
\hline
Частота вращения $n$ & \varN & об/мин \\
\hline
Вязкость смазки $\eta$ & \varEta & Па$\cdot$с \\
\hline
Диапазон эксцентриситета $\varepsilon$ & 0{,}05-0{,}80 & -- \\
\hline
Размер сетки & $500 \times 500$ & -- \\
\hline
Параметр релаксации $\omega$ & 1{,}5 & -- \\
\hline
Критерий сходимости $\delta$ & $< 10^{-5}$ & -- \\
\hline
Макс. число итераций & 50000 & -- \\
\hline
\end{tabularx}
\end{table}

Линейная скорость поверхности вала при заданной частоте вращения составляет $U = \pi D n / 60 = \pi \cdot 0{,}070 \cdot 2980 / 60 \approx 10{,}91$~м/с. Расчёт выполнен для 15 значений эксцентриситета в указанном диапазоне.

\subsection{Распределение зазора и давления}

Распределение безразмерного давления $P(\varphi)$ в среднем сечении ($Z = 0$) при $\varepsilon = 0{,}6$ для гладкого и текстурированного подшипников представлено на рисунке~\ref{fig:pressure_2d}.

\begin{figure}[H]
\centering
\includegraphics[width=\textwidth]{figures/pressure_2d.png}
\caption{Распределение давления $P(\varphi)$ при $Z=0$ для гладкого и текстурированного подшипников ($\varepsilon = 0{,}6$)}
\label{fig:pressure_2d}
\end{figure}

Для гладкого подшипника характерно наличие единственного максимума давления в области минимального зазора ($\varphi \approx \pi$). Давление нарастает в сужающейся части зазора (от $\varphi = 0$ до $\varphi \approx \pi$) и обращается в нуль в расширяющейся части вследствие кавитации. Область положительного давления занимает примерно половину окружности. При наличии конического микрорельефа на кривой давления появляются плавные локальные возмущения, соответствующие углублениям. В отличие от профилей с резкой границей, конический профиль с линейно убывающей глубиной создаёт более мягкий клиновый зазор, поэтому возмущения давления имеют сглаженный вид без выраженных пиков и провалов.

На рисунке~\ref{fig:clearance_2d} показан профиль зазора $H(\varphi)$ в среднем сечении при $\varepsilon = 0{,}6$.

\begin{figure}[H]
\centering
\includegraphics[width=\textwidth]{figures/clearance_2d.png}
\caption{Распределение зазора $H(\varphi)$ при $Z=0$ для подшипника с коническими углублениями ($\varepsilon = 0{,}6$)}
\label{fig:clearance_2d}
\end{figure}

На фоне плавного изменения зазора по закону $H = 1 + \varepsilon\cos\varphi$ наблюдаются локальные увеличения, соответствующие коническим углублениям. Конический профиль характеризуется максимальной глубиной в центре углубления и линейно убывающей глубиной к периферии: $\Delta h = h_p(1 - r)$ при $r \leq 1$. Максимальная глубина каждого углубления составляет $\Delta H_{\max} = h_p/c = 0{,}2$ в безразмерных единицах, на границе углубления глубина обращается в нуль.

Трёхмерные распределения зазора и давления при $\varepsilon = 0{,}6$ представлены на рисунках~\ref{fig:clearance_3d} и~\ref{fig:pressure_3d}.

\begin{figure}[H]
\centering
\includegraphics[width=\textwidth]{figures/clearance_3d.png}
\caption{Трёхмерное распределение зазора $H(\varphi, Z)$ ($\varepsilon = 0{,}6$)}
\label{fig:clearance_3d}
\end{figure}

На рисунке~\ref{fig:clearance_3d} видна регулярная структура микрорельефа: 8 конических углублений по окружности и 11 рядов по длине подшипника в зоне $90^\circ\text{ - }270^\circ$. Каждое углубление имеет вид перевёрнутого конуса с линейным изменением глубины от центра к границе.

\begin{figure}[H]
\centering
\includegraphics[width=\textwidth]{figures/pressure_3d.png}
\caption{Трёхмерное распределение давления $P(\varphi, Z)$ ($\varepsilon = 0{,}6$)}
\label{fig:pressure_3d}
\end{figure}

На рисунке~\ref{fig:pressure_3d} помимо основного максимума давления, обусловленного макрогеометрией клинового зазора, наблюдаются плавные локальные возмущения давления над каждым коническим углублением. Линейное изменение глубины формирует менее выраженный микроклин, чем профили с резкой границей, поэтому локальные пики и провалы давления оказываются сглаженными. Из-за кавитационного ограничения ($P \geq 0$) отрицательная часть возмущения обрезается, и суммарный вклад каждого углубления остаётся положительным, хотя по абсолютной величине он меньше, чем у углублений с более резкой геометрией. Конический профиль обеспечивает плавное взаимодействие смазочного слоя с микрорельефом, что благоприятно сказывается на стабильности течения, но ограничивает прирост несущей способности.

Анализ трёхмерных распределений позволяет отметить следующие особенности. Давление распределено неравномерно по длине подшипника -- максимум достигается в среднем сечении ($Z = 0$), а к торцам ($Z = \pm 1$) давление падает до нуля вследствие торцевых утечек. Влияние микрорельефа наиболее заметно вблизи среднего сечения, где базовое давление максимально. На периферийных рядах углублений (вблизи торцов) эффект текстурирования ослабевает из-за малого уровня фонового давления.

\subsection{Сравнительный анализ интегральных показателей}

Интегральные характеристики гладкого подшипника и подшипника с коническими углублениями для выбранных значений эксцентриситета представлены в таблицах~\ref{tab:smooth_results} и~\ref{tab:textured_results}.

\begin{table}[H]
\small
\centering
\caption{Характеристики гладкого подшипника}
\label{tab:smooth_results}
\begin{tabular}{|c|c|c|c|}
\hline
$\boldsymbol{\varepsilon}$ & $\boldsymbol{F}$\textbf{, Н} & $\boldsymbol{\mu}$ & $\boldsymbol{Q}$\textbf{, л/с} \\
\hline
0{,}05 & 276{,}2 & 0{,}2153 & 0{,}2398 \\
\hline
0{,}16 & 888{,}9 & 0{,}0678 & 0{,}2395 \\
\hline
0{,}26 & 1571{,}1 & 0{,}0395 & 0{,}2388 \\
\hline
0{,}37 & 2389{,}0 & 0{,}0272 & 0{,}2379 \\
\hline
0{,}48 & 3452{,}2 & 0{,}0201 & 0{,}2368 \\
\hline
0{,}59 & 4975{,}6 & 0{,}0153 & 0{,}2357 \\
\hline
0{,}69 & 7472{,}8 & 0{,}0117 & 0{,}2347 \\
\hline
0{,}80 & 12572{,}9 & 0{,}0085 & 0{,}2338 \\
\hline
\end{tabular}
\end{table}

\begin{table}[H]
\small
\centering
\caption{Характеристики подшипника с коническими углублениями}
\label{tab:textured_results}
\begin{tabular}{|c|c|c|c|}
\hline
$\boldsymbol{\varepsilon}$ & $\boldsymbol{F}$\textbf{, Н} & $\boldsymbol{\mu}$ & $\boldsymbol{Q}$\textbf{, л/с} \\
\hline
0{,}05 & 289{,}3 & 0{,}2043 & 0{,}2415 \\
\hline
0{,}16 & 907{,}1 & 0{,}0660 & 0{,}2412 \\
\hline
0{,}26 & 1596{,}9 & 0{,}0386 & 0{,}2405 \\
\hline
0{,}37 & 2427{,}1 & 0{,}0266 & 0{,}2396 \\
\hline
0{,}48 & 3516{,}4 & 0{,}0196 & 0{,}2385 \\
\hline
0{,}59 & 5110{,}1 & 0{,}0147 & 0{,}2374 \\
\hline
0{,}69 & 7857{,}7 & 0{,}0109 & 0{,}2363 \\
\hline
0{,}80 & 14195{,}2 & 0{,}0074 & 0{,}2353 \\
\hline
\end{tabular}
\end{table}

Нагрузочная способность $F$ нелинейно возрастает с увеличением эксцентриситета для обоих вариантов подшипника. Текстурированный подшипник обладает большей нагрузочной способностью при всех значениях $\varepsilon$, однако разница между ним и гладким аналогом сравнительно невелика. При высоких эксцентриситетах ($\varepsilon = 0{,}80$) эффект текстурирования выражен наиболее ярко: нагрузочная способность текстурированного подшипника (14195{,}2~Н) превышает гладкий аналог (12572{,}9~Н) на 12{,}9\%.

Коэффициент трения $\mu$ монотонно убывает при увеличении $\varepsilon$ и для текстурированного подшипника ниже при всех значениях эксцентриситета. Расход смазки $Q$ слабо зависит от $\varepsilon$ и у текстурированного подшипника несколько выше за счёт увеличения эффективного объёма смазочного слоя в зоне углублений.

Для количественной оценки эффективности текстурирования в таблице~\ref{tab:comparison} представлены результаты сравнения при $\varepsilon = 0{,}6$.

\begin{table}[H]
\small
\centering
\caption{Сравнение характеристик при $\varepsilon = 0{,}6$}
\label{tab:comparison}
\begin{tabular}{|l|c|c|c|}
\hline
\textbf{Характеристика} & \textbf{Гладкий} & \textbf{С углублениями} & $\boldsymbol{\Delta}$\textbf{, \%} \\
\hline
Нагрузка $F$, Н & 5234{,}7 & 5386{,}6 & $+2{,}90$ \\
\hline
Коэфф. трения $\mu$ & 0{,}01480 & 0{,}01420 & $-4{,}06$ \\
\hline
Расход $Q$, л/с & 0{,}2356 & 0{,}2373 & $+0{,}71$ \\
\hline
\end{tabular}
\end{table}

Нанесение конических углублений увеличивает несущую способность подшипника на 2{,}90\% за счёт микрогидродинамического эффекта. Прирост невелик по сравнению с профилями, имеющими более выраженную клиновую геометрию, что объясняется плавным линейным изменением глубины конуса: на входной кромке углубления формируется лишь слабый дополнительный микроклин. Коэффициент трения снижается на 4{,}06\% вследствие увеличения средней толщины смазочного слоя в зоне углублений. Расход смазки возрастает на 0{,}71\%, что незначительно сказывается на требованиях к системе маслоснабжения.

На рисунках~\ref{fig:load_vs_eps} - \ref{fig:flow_vs_eps} представлены зависимости интегральных характеристик от эксцентриситета для обоих вариантов подшипника.

\begin{figure}[H]
\centering
\includegraphics[width=\textwidth]{figures/load_vs_epsilon.png}
\caption{Нагрузочная способность $F(\varepsilon)$ для гладкого и текстурированного подшипников}
\label{fig:load_vs_eps}
\end{figure}

Кривая $F(\varepsilon)$ для текстурированного подшипника проходит выше во всём диапазоне эксцентриситетов, но кривые расположены близко друг к другу. С ростом $\varepsilon$ разрыв между ними постепенно увеличивается: если при $\varepsilon = 0{,}05$ разница составляет около 13~Н, то при $\varepsilon = 0{,}80$ она достигает примерно 1620~Н. Нелинейный характер зависимости обусловлен степенной зависимостью давления от толщины зазора ($P \sim h^{-3}$), однако для конического профиля нелинейность выражена слабее, чем для углублений с резкими кромками.

\begin{figure}[H]
\centering
\includegraphics[width=\textwidth]{figures/friction_vs_epsilon.png}
\caption{Коэффициент трения $\mu(\varepsilon)$ для гладкого и текстурированного подшипников}
\label{fig:friction_vs_eps}
\end{figure}

Коэффициент трения текстурированного подшипника ниже гладкого при всех значениях $\varepsilon$. Обе кривые монотонно убывают, асимптотически приближаясь к нулю при $\varepsilon \to 1$. Абсолютная разница между кривыми составляет около $\Delta\mu = 0{,}011$ при $\varepsilon = 0{,}05$ и $\Delta\mu = 0{,}001$ при $\varepsilon = 0{,}80$. В относительном выражении снижение трения сопоставимо во всём диапазоне эксцентриситетов и находится на уровне нескольких процентов.

\begin{figure}[H]
\centering
\includegraphics[width=\textwidth]{figures/flow_vs_epsilon.png}
\caption{Расход смазки $Q(\varepsilon)$ для гладкого и текстурированного подшипников}
\label{fig:flow_vs_eps}
\end{figure}

Расход смазки текстурированного подшипника незначительно превышает расход гладкого. Разница составляет порядка 0{,}7\% и практически не зависит от эксцентриситета.

Таким образом, конический микрорельеф обеспечивает одновременное увеличение нагрузочной способности и снижение коэффициента трения при незначительном росте расхода. Эффект усиливается при увеличении эксцентриситета, однако в абсолютном и относительном выражении он скромнее, чем у профилей с более резкой клиновой геометрией.

\subsection*{Выводы по главе 3}
\phantomsection
\addcontentsline{toc}{subsection}{Выводы по главе 3}

Выполнен расчёт характеристик радиального подшипника скольжения с коническим микрорельефом для 15 значений эксцентриситета. Получены поля давления и зазора, построены зависимости нагрузочной способности, коэффициента трения и расхода смазки от эксцентриситета. Сравнительный анализ при $\varepsilon = 0{,}6$ показал увеличение нагрузочной способности на $2{,}90$\%, снижение коэффициента трения на $4{,}06$\% и рост расхода на $0{,}71$\%, что подтверждает положительное, хотя и умеренное, влияние конического текстурирования.
